\documentclass[11pt]{article}
\usepackage[T1]{fontenc}
\usepackage[utf8]{inputenc}
\usepackage{lmodern,microtype,geometry,amsmath,amssymb,amsthm,mathtools}
\usepackage[colorlinks=true,linkcolor=blue,urlcolor=blue,citecolor=blue]{hyperref}
\geometry{margin=1in}
\setlength{\emergencystretch}{2em}
\newtheorem{theorem}{Theorem}[section]
\newtheorem{lemma}[theorem]{Lemma}
\newtheorem{proposition}[theorem]{Proposition}
\newtheorem{corollary}[theorem]{Corollary}
\theoremstyle{remark}\newtheorem{remark}[theorem]{Remark}
\newcommand{\R}{\mathbb R}\newcommand{\E}{\mathbb E}
\newcommand{\norm}[1]{\left\lVert#1\right\rVert}
\newcommand{\cX}{\mathcal X}\newcommand{\cS}{\mathcal S}
\newcommand{\clip}{\operatorname{clip}}
\hypersetup{pdftitle={Polynomial-depth mistake converses for neural priors with controlled internal amplification},pdfauthor={ChatGPT-6 Astra; Claude Fable 5 and Opus 5; Kaarel Hanni; Dmitry Vaintrob}}
\title{Polynomial-depth mistake converses for neural priors\\with controlled internal amplification}
\author{ChatGPT-6 Astra \and Claude Fable 5 and Opus 5 \and Kaarel H\"anni \and Dmitry Vaintrob\thanks{Authors are listed in alphabetical order.}}
\date{Research draft, 22 October 2026}
\begin{document}\maketitle
\begin{abstract}
Does making few online prediction mistakes force a Bayesian neural-network
learner's target to have a small implementation, with polynomial rather
than exponential dependence on depth? We prove such a converse for a
Gaussian tanh prior restricted by bounded output magnitude and a
label-independent certificate controlling propagation of internal noise.
The certificate allows compensated layer expansion and includes networks
whose internal matrices have operator norm at most one. It is preserved by
cloning and direct sums. An exact singleton correction then increases the
square root of minimum energy by only a polynomial-depth additive cost.
A deterministic robustness compressor turns the resulting quadratic
energy bound into an explicit circuit-width bound with fourth-power
mistake dependence and no exponent depending on depth. The theorem
concerns the limit taking prior variance to zero at each fixed width and
then letting width tend to infinity. It handles tied minimizers,
posterior sampling, majority prediction, and posterior-confidence
failures, in any complete input-reveal order. The original fixed-width
converse, and cheap implementations of every small deep Boolean circuit
under this stronger restriction, remain open. These results are not
Lean verified.
\end{abstract}

\section{Question and scope}
An online learner receives an input, predicts its Boolean label, and then
receives the true label. Every label is revealed, including after correct
predictions. A converse asks whether a small total mistake count implies
that the entire target has a small circuit. It requires sufficient input
coverage: observations on a small easy subset cannot constrain the target
elsewhere.

For minimum-size Boolean-circuit induction, an incorrect prediction can
be repaired by adding a short exception circuit. This immediately bounds
circuit size by the mistake count. The analogous neural argument has two
difficulties. A large wrong real-valued output may require an expensive
correction, and a low-energy neural network need not have a known
polynomial-depth compressor. We address these separately: a fixed output
cap makes corrections inexpensive, while a stable internal-noise
certificate supplies a polynomial-depth compressor.

Our conclusion concerns bounded-weight tanh-network width at the stated
depth, rather than unrestricted-depth NAND gate count. The new argument
combines an exact singleton construction with the deterministic compressor
underlying the earlier polynomial-depth learning result
\cite{quartic}. We reproduce the required compression argument below.
The prior restriction is specified in weights and outputs, without
reference to training labels, target complexity, or mistake count.

The order of limits is substantive. We take the strong-prior limit at
each fixed width first, and only then let width tend to infinity. This
supplies the extra neurons used by an exception branch. We do not assert
that the converse holds at every fixed realizable width, or that this
limit is one of the finite-exponent width-first posteriors.

\section{A prior with controlled internal amplification}
Fix $d\ge1$, $L\ge2$, $H=L-1$, $U>1$, and $\Lambda>1$.
Let $\cX=\{-1,1\}^d$ and $\bar x=(x,1)$. A network with $H$ hidden
layers of common width $N$ has
\[
 h^1=\tanh(W^1\bar x),\qquad
 h^s=\tanh(W^sh^{s-1})\quad(2\le s\le H),\qquad
 z_w=v^\top h^H.
\]
Its normalized Gaussian energy and readout energy are
\[
 \kappa_N(w)=\frac dN\norm{W^1}_F^2+
 \sum_{s=2}^H\norm{W^s}_F^2+N\norm v_2^2,
 \qquad E_v(w)=N\norm v_2^2.
\]
The independent raw Gaussian coordinates have variances
$N\sigma^2/d$ in the first layer, $\sigma^2$ internally, and
$\sigma^2/N$ in the readout. The density is proportional to
$\exp[-\kappa_N/(2\sigma^2)]$.

For diagonal matrices $D_2,\ldots,D_H$ with entries in $[0,1]$, put
\[
 B_j(w,D)=v^\top D_HW^H\cdots D_{j+1}W^{j+1}
 \quad(1\le j<H),\qquad B_H(w,D)=v^\top.
\]
Define the certified squared hidden-noise gain
\begin{equation}\label{eq:gain}
 G_N(w)^2=N\max_{D_2,\ldots,D_H}
                     \sum_{j=1}^H\norm{B_j(w,D)}_2^2.
\end{equation}
The maximum is over a compact finite-dimensional cube, so it exists and
is continuous in the weights. For $H=1$, it is simply $E_v$.

Let $\cS_{N,\Lambda,U}$ be the following open parameter set:
\begin{equation}\label{eq:support}
 \frac12<|z_w(x)|<U\quad(x\in\cX),\qquad
 G_N(w)^2<\Lambda^2H E_v(w).
\end{equation}
The prior is the raw Gaussian conditioned on this set. Training further
conditions on exact agreement of output signs with the revealed labels.
Thus we retain the original global lower-margin convention and add an
upper output cap and an internal-amplification certificate.

\begin{remark}[Interpretation]
The diagonals in \eqref{eq:gain} range over all possible slopes of tanh,
not only slopes at the clean forward pass. The condition therefore
certifies finite internal disturbances rather than only an infinitesimal
Jacobian at one state. Requiring $\norm{W^s}_{\rm op}\le1$ at every
internal layer is sufficient, since then $G_N^2\le H E_v$.
The certificate is more flexible: expansion in a layer is allowed when
subsequent contractions or readout directions control its effect.
It is stronger than checking only the perturbations encountered on a
particular data set.
\end{remark}

\begin{lemma}[Certified finite-noise stability]\label{lem:noise}
Add disturbances after every hidden activation and at the readout, and
measure them by
\[
 \norm\eta_*^2=\sum_{s=1}^H\norm{\eta^s}_2^2/N+|\eta^{\rm out}|^2.
\]
Every $w$ satisfying the gain condition in \eqref{eq:support} obeys
\[
 |z_w(x;\eta)-z_w(x)|
 \le\Lambda\sqrt{1+HE_v(w)}\,\norm\eta_*.
\]
In particular, with $c=1/(16\Lambda)$ and
$r_c(w)=c/\sqrt{1+HE_v(w)}$, its clipped output
$T(z)=\clip(z,-1/2,1/2)$ changes by at most $1/16$ whenever
$\norm\eta_*\le r_c(w)$.
\end{lemma}
\begin{proof}
The difference of two tanh values is their input difference times a
number in $[0,1]$. Thus the hidden-state differences satisfy
$\Delta h^1=\eta^1$ and
$\Delta h^s=D_sW^s\Delta h^{s-1}+\eta^s$ for suitable diagonals
$D_s$. Expanding this recursion gives
$\Delta z=\sum_jB_j(w,D)\eta^j+\eta^{\rm out}$.
Cauchy--Schwarz bounds its magnitude by
$\sqrt{1+G_N^2}\norm\eta_*$, and
$1+G_N^2\le\Lambda^2(1+HE_v)$ since $\Lambda>1$.
Clipping is one-Lipschitz.
\end{proof}

\section{Statement of the converse}
For a Boolean function $g$, let $a_N(g)$ be the infimum of
$\kappa_N(w)$ over $w\in\cS_{N,\Lambda,U}$ with sign function $g$,
with value $+\infty$ when there is no such implementation. Set
$a_\infty(g)=\inf_N a_N(g)$.
A minimum-action predictive rule is one supported, at each history, on
functions minimizing $a_\infty$ among those agreeing with that history.
It may sample such a function, choose one deterministically, or predict
by the majority of such a distribution. Confidence failure means that
the correct label has probability at most a fixed $\theta\in(0,1)$.

Let $C_L(f)$ be the smallest integer $C$ such that a depth-$L$, width-$C$
tanh network has $f(x)z(x)>1/2$ everywhere and every weight has absolute
value at most $C$. The network witnessing this conclusion need not
satisfy the additional prior restriction.

Define explicit budgets
\begin{align*}
 A_d&=\tfrac12\log(4(d+1)),& F_d&=1+A_d^2d(d+1),\\
 D_{d,H}&=dF_d+H,& Q_H&=4(H+65),\\
 P_{d,H,U}&=D_{d,H}+Q_H(U+1)^2.
\end{align*}

\begin{theorem}[Polynomial-depth mistake converse]\label{thm:main}
Let $f:\cX\to\{-1,1\}$ and reveal every point of $\cX$, in any order.
For any minimum-action predictive rule as above, let $m$ be either its
actual total mistake count on a complete-revelation path or its expected
total mistake count, using any one of the stated prediction conventions.
Put
\[
 K=1+P_{d,H,U}(m+1)^2.
\]
Then
\begin{equation}\label{eq:main}
 C_L(f)\le
 \left\lceil65536\Lambda^2K(1+HK)
                   [d+\log(400L)]\right\rceil.
\end{equation}
This is a polynomial bound jointly in $d,L,U,\Lambda,m$, with
logarithmic factors displayed and with leading fourth-power dependence
on $m$. No exponent or implicit constant depends on depth.

Every subsequential predictive limit obtained by first taking
$\sigma\downarrow0$ at fixed $N$ for the prior
\eqref{eq:support}, and then taking $N\to\infty$, is a
minimum-action rule and satisfies the theorem. The same conclusion holds
on an infinite iid stream with full support, since every domain point is
eventually revealed and repeated points cause no further mistakes.
\end{theorem}

The theorem also holds if the gain certificate is replaced throughout by
$\norm{W^s}_{\rm op}<1$ for every internal layer. This simpler support
has its own minimum energies and predictive limits; its closure and
singleton constructions below are unchanged. Its networks satisfy
Lemma~\ref{lem:noise} with any fixed $\Lambda>1$, for example $\Lambda=2$.

The statement is a converse. It does not establish that every small
Boolean circuit has a low-energy implementation in this support.
A fixed output cap controls a true amplitude obstruction in the repair
argument; it is not inferred from small sign error or bounded hidden
activations. A hard-conditioning prior need not describe ordinary SGD.

\section{Cloning, direct sums, and the Gaussian limit}

\begin{lemma}[Closure under cloning and direct sums]\label{lem:closure}
Uniformly clone each neuron of a width-$n$ network $k$ times, divide its
outgoing weights by $k$, and pad with zero neurons to any width $N\ge kn$.
Its function is unchanged. If $\lambda=N/(kn)$, then
\[
 E'_v=\lambda E_v,\qquad G_N(w')^2=\lambda G_n(w)^2.
\]
The first-layer energy is divided by $\lambda$, internal energy is
unchanged, and readout energy is multiplied by $\lambda$.
Thus gain support and output caps are preserved.

For a block-diagonal direct sum of widths $n_1,n_2$, with outputs added
and ambient width $N=n_1+n_2$,
\[
 E_v=\sum_{i=1}^2\frac N{n_i}E_{v_i},\qquad
 G_N^2=\sum_{i=1}^2\frac N{n_i}G_{n_i}^2.
\]
Consequently the strict gain certificate is preserved. The output cap
is preserved whenever the resulting real output satisfies it.
\end{lemma}
\begin{proof}
On active clone coordinates the internal matrix is
$W^s\otimes(J_k/k)$ and the readout is $v\otimes\mathbf1_k/k$.
Every path row $B_j$ is the old path row, with diagonal entries replaced
by their within-clone averages, tensored with $\mathbf1_k^\top/k$.
Its squared norm is divided by $k$. All averaged diagonals in $[0,1]$
are attainable by taking identical slopes in each clone group, so the
supremum is exactly the old one divided by $k$. Padding has no effect
on these path rows. This proves the gain identity; direct calculation
gives the remaining cloning energies.

For a direct sum, every path row splits into the two blocks. The slope
choices in the blocks are independent, so the supremum of the sum of
squared norms is the sum of the two suprema. This proves the identities.
For the operator-norm support, cloning preserves internal operator norms
and a block-diagonal matrix has norm equal to the maximum of its blocks.
\end{proof}

For an exactly realizable real function $u$, let $E(u)$ be the infimum
energy over all widths satisfying the gain certificate, without a lower
margin or upper output constraint. Zero functions may be assigned energy
zero by a limit. Lemma~\ref{lem:closure} gives
\begin{equation}\label{eq:triangle}
 \sqrt{E(u+v)}\le\sqrt{E(u)}+\sqrt{E(v)}.
\end{equation}
Indeed, clone two fixed branches into asymptotic proportions $1-\alpha$
and $\alpha$. Their first-layer energies are multiplied by these
proportions, internal energies do not change, and readout energies are
divided by them. The total is at most
$E(u)/(1-\alpha)+E(v)/\alpha$, after taking infima and approximation
limits. Optimizing $\alpha$ proves \eqref{eq:triangle}. The same
argument works for the operator-norm support.

\begin{proposition}[Connection with the stated Gaussian learner]
\label{prop:limits}
Every Boolean function has finite $a_\infty$, and
$a_N(g)\to a_\infty(g)$ as $N\to\infty$. The open-support Gaussian
prior's fixed-width strong-prior posterior is supported in the limit on
functions of smallest $a_N$ consistent with the history. Every subsequent
width-limit accumulation point is supported on smallest $a_\infty$.
\end{proposition}
\begin{proof}
Finite realizability follows from the exact singleton construction in the
next section: at one hidden layer sum the singleton expansions with
coefficients $g(x)$ to obtain output $g$ exactly. At greater depth,
combine the first-layer features with sufficiently small coefficients to
make the second-layer operator norm less than one, obtaining a scalar
hidden state of equal magnitude and sign $g(x)$. Contractive repeated
copies preserve that sign and common magnitude; the readout restores
output exactly to $g$. Cloning the feature rows with suitable proportions
makes the second-layer norm arbitrarily close to its variation-norm bound.
Thus all internal operator norms are strictly below one and the output
is exactly $\pm1$, within the fixed cap. A sufficiently large finite
width suffices.

Cloning a fixed width-$n$ realization with $k=\lfloor N/n\rfloor$
shows that $\limsup_N a_N(g)\le a_n(g)$, because
$N/(kn)\to1$. Taking the infimum over $n$ and using the reverse
inequality from the infimum definition proves convergence.

At fixed width, the event for any given function is open in parameter
space. In weighted coordinates the raw Gaussian is $N(0,\sigma^2I)$.
The elementary Gaussian rate calculation yields
$2\sigma^2\log\mu_{N,\sigma}(U_g)\to-a_N(g)$: integrate over a
small ball for the lower bound, and use the Gaussian norm tail for the
upper bound. The support-normalization factor cancels in ratios.
There are only finitely many functions, so strictly higher-energy
functions disappear after conditioning on a noiseless history.
The same finite-set argument and convergence of $a_N$ prove the final
claim. Simultaneous subsequences for all finite distinct-input histories
exist because their number is finite.
\end{proof}

\section{Exact singleton corrections with controlled amplification}

\begin{lemma}[An inexpensive exact point indicator]\label{lem:singleton}
For every $a\in\cX$ and $t\in\R$, the function
$tI_a(x)=t\mathbf1\{x=a\}$ has implementations at depth $L$ with
internal operator norms strictly below one and with infimum energy at most
\[
 E(tI_a)\le D_{d,H}+Q_Ht^2.
\]
\end{lemma}
\begin{proof}
There are $d+1$ affine-tanh features $h_0,\ldots,h_d$ satisfying
\begin{equation}\label{eq:singleton-expansion}
 I_a=\sum_{j=0}^d b_jh_j,\qquad \sum_j|b_j|\le4,
 \qquad \norm{\text{first-layer row }j}_2^2\le F_d.
\end{equation}
To see this, take $h_0=\tanh(1)$ and
\[
 h_j(x)=\tanh\left(A_d[a\cdot x-(d-2j+1)]\right),\quad 1\le j\le d.
\]
On the shells $s=(d-a\cdot x)/2\in\{0,\ldots,d\}$, the columns
$1,(1+h_1)/2,\ldots,(1+h_d)/2$ are entrywise within
$e^{-2A_d}$ of the staircase columns
$1,\mathbf1\{s<1\},\ldots,\mathbf1\{s<d\}$.
The staircase inverse takes differences and a final endpoint, so its
induced $\ell^1$ norm is at most two. The perturbation in inverse
coordinates has norm at most
$2(d+1)e^{-2A_d}=1/2$. The target singleton has staircase coefficient
norm one; a Neumann-series inverse gives coefficients of total magnitude
at most two. Expanding the affine columns and using
$1=h_0/\tanh(1)$ increases that norm by a factor less than two.
The row bound follows from $d$ coefficients of magnitude $A_d$ and a
bias of magnitude at most $A_dd$. This proves
\eqref{eq:singleton-expansion}.

Write $V=\sum_j|b_j|\le4$ and omit zero coefficients. At large width
clone feature $j$ in proportion $p_j=|b_j|/V$. If $H=1$, split its
readout coefficient $tb_j$ equally among its copies. The limiting
readout energy is $t^2\sum_jb_j^2/p_j=t^2V^2\le16t^2$,
and first-layer energy is at most $dF_d$.

If $H\ge2$, let each second-layer neuron receive $b_j/8$ times feature
$j$, distributed across its copies. The resulting state is
$a_2I_a$ with $a_2=\tanh(1/8)$. The second-layer squared operator
norm and Frobenius energy both tend to
$V^2/64\le1/4$, so they are below one at all sufficiently accurate
finite allocations. Later hidden matrices are constant matrices
$\beta J/N$, where $\beta=1-1/(2H)<1$. Thus every internal matrix
is strictly contractive, and the state is $a_sI_a$, where
$a_s=\tanh(\beta a_{s-1})$. Readout entries are $t/(Na_H)$.

For $y\ge0$, $\tanh y\ge y/\sqrt{1+y^2}$, because
$\operatorname{arsinh}y\le y$ and
$\tanh(\operatorname{arsinh}y)=y/\sqrt{1+y^2}$.
Consequently
\[
 a_2^{-2}\le65,\qquad
 a_s^{-2}\le1+\beta^{-2}a_{s-1}^{-2},\qquad
 a_H^{-2}\le\beta^{-2H}(65+H)\le4(65+H)=Q_H.
\]
The last inequality uses
$-\log(1-1/(2H))\le1/(2H-1)$ and $H\ge2$.
The first-layer cost is at most $dF_d$; the total internal cost is
at most $H$ in the limiting construction; readout energy is at most
$Q_Ht^2$. Every finite allocation computes the singleton exactly,
so no approximation consumes any pre-existing margin slack.
\end{proof}

A constant output $+1$ or $-1$ also has a strictly contractive
implementation of energy at most $d+5H+8$. Use identical first-layer
features $\tanh(1)$, the same contractive relays, and a readout restoring
unit magnitude. The inverse-square calculation above bounds the required
readout energy by $4(H+2)$. In particular, the initial minimum action is
at most $P_{d,H,U}$.

\section{Mistakes bound the minimum action}
Let $K_t$ be the smallest $a_\infty$ among functions consistent with the
first $t$ revealed labels. This minimum exists, since the domain has
finitely many functions and all have finite action. It is nondecreasing.
If $K_t>K_{t-1}$, every function minimizing the previous action has the
wrong label at the new input. Such a round is a forced mistake for every
prediction convention in Theorem~\ref{thm:main}. This remains true with
arbitrary mixtures of tied minimizers.

Choose a minimizing function before such an increase and a capped
real-output realization $z$ with energy arbitrarily close to its action.
If the next input is $a$ with label $y$, replace its output by
\[
 z_{\rm new}(x)=z(x)+(y-z(a))I_a(x).
\]
Every other output is unchanged, and the new value at $a$ is exactly
$y\in\{-1,1\}$. All previous labels, the strict lower margin, and the
upper cap remain valid. Since $|z(a)|<U$, the singleton amplitude is at
most $U+1$. Lemmas~\ref{lem:closure} and \ref{lem:singleton}, followed
by approximation of the energy infima, give
\[
 \sqrt{K_t}\le\sqrt{K_{t-1}}+\sqrt{P_{d,H,U}}
\]
at an increase step. At other steps $K_t=K_{t-1}$. If $B$ is the number
of increase steps after full revelation, then
\begin{equation}\label{eq:energy-mistakes}
 a_\infty(f)\le P_{d,H,U}(B+1)^2
                    \le P_{d,H,U}(M+1)^2,
\end{equation}
where $M$ is total mistakes under any stated convention.

For random orders or random predictions, the same conclusion holds with
$m=\E M$. Invert \eqref{eq:energy-mistakes} first:
$M\ge\sqrt{a_\infty(f)/P_{d,H,U}}-1$ on every complete path.
The right side is deterministic, so expectation and rearrangement give
$a_\infty(f)\le P_{d,H,U}(m+1)^2$. This does not reverse Jensen's
inequality. Predictable adaptive input orders are allowed, provided all
points are eventually revealed.

\section{Polynomial-depth reconstruction from energy}
For completeness we give the deterministic compression statement needed
to turn \eqref{eq:energy-mistakes} into circuit complexity. It is the
pruning-and-Bernstein construction of \cite{quartic}, used without its
Gaussian or statistical parts.

\begin{lemma}[Robust energy compression]\label{lem:compression}
Suppose a width-$N$, depth-$L$ tanh network has $\kappa_N(w)\le K$,
where $K\ge1$, and its clipped outputs change by at most $1/16$ under
all disturbances of norm at most $r$, where $0<r\le1/6$. Put
\[
 m_*=\left\lceil256Kr^{-2}[d+\log(400L)]\right\rceil.
\]
There is an actual width-$m_*$, depth-$L$ tanh network with energy at
most $4K$, weights bounded in magnitude by $2\sqrt{m_*K}$, and
clipped outputs differing from the original's by at most $1/16$ on the
entire Boolean cube.
\end{lemma}
\begin{proof}
Write $k_s=\norm{W^s}_F^2$ for internal layers and
$k_L=N\norm v_2^2$, so their sum is at most $K$.
In each internal row and the readout, delete entries of magnitude above
$8\sqrt K\norm{\text{row}}_2/(r\sqrt N)$; zero rows need no change.
The absolute removed row mass is at most its squared norm divided by
this threshold. On a vector in $[-1,1]^N$, pruning therefore has
normalized stage error at most $r\sqrt{k_s/K}/8$.

At each hidden layer independently sample $m_*$ uniform source indices
with replacement. Keep sampled first-layer rows, and multiply subsequent
sampled source weights by $N/m_*$. Work first with a virtual pass retaining
all destination rows. Its clean virtual hidden states remain in
$[-1,1]^N$ and are independent of the fresh source indices at the next
stage. For one uniform source index $I$, the internal vector summand is
$Y=N\overline W^s_{:,I}\widetilde h_I$. It obeys
\[
 \E\norm{Y-\E Y}_2^2\le Nk_s,\qquad
 \norm{Y-\E Y}_2\le16\sqrt{NKk_s}/r.
\]
We use the Hilbert-space Bernstein inequality
\cite[Appendix A, Corollary A.2]{tarres-yao}: independent centered
Hilbert-space variables $X_i$ with $\norm{X_i}\le b$ and
$\sum_i\E\norm{X_i}^2\le V$ satisfy
\[
 \Pr\!\left\{\left\|\sum_iX_i\right\|\ge t\right\}
 \le2\exp\!\left[-\frac{t^2}{2(V+bt/3)}\right].
\]
For the sample average at tolerance
$r\sqrt{Nk_s/K}/8$, this bounds the sampling failure probability by
$2\exp[-m_*r^2/(256K)]$. The readout calculation is the same without
factors of $\sqrt N$. A union bound over at most $L$ stages and $2^d$
inputs is less than $1/200$ by the chosen budget.

On success the pruning and sampling errors combine to give normalized
stage errors at most $r\sqrt{k_s/K}/4$. Define postactivation disturbances
as the difference between each virtual state and tanh applied to the
original weights on the preceding virtual state. One-Lipschitzness makes
the total disturbance norm at most $r/4$, uniformly on the cube. A
readout disturbance makes the virtual output exact. Thus its clipped
output differs from the original by at most $1/16$.

The stored sampled coordinates form an ordinary width-$m_*$ network.
Independent index lists in adjacent layers show that its expected
normalized energy is at most $K$: the first-layer contribution is
unchanged in expectation, and every remaining contribution is at most
its original value. Hence its energy is at most $4K$ with probability
at least $3/4$. The energy and all-input-noise events intersect. Choose
one such network. The energy bound gives every coordinate magnitude at
most $2\sqrt{m_*K}$, including the first layer since $d\ge1$.
\end{proof}

To finish Theorem~\ref{thm:main}, choose an admissible realization of $f$
with energy below $K=1+P_{d,H,U}(m+1)^2$. Lemma~\ref{lem:noise} applies
at radius
\[
 r=\frac1{16\Lambda\sqrt{1+HK}},
\]
which is at most $1/6$ and no greater than the network's own certified
radius. Lemma~\ref{lem:compression} gives precisely the width in
\eqref{eq:main}. Since the original output has signed margin above
$1/2$, its signed clipped output is exactly $1/2$. The compressed output
has signed margin at least $7/16$. Multiply its readout by two to obtain
margin at least $7/8$, strictly above $1/2$. All weights are then at most
$4\sqrt{m_*K}\le m_*$, because the displayed budget has $m_*\ge16K$.
This proves $C_L(f)\le m_*$.

For fixed $U,\Lambda$, the order of the bound is, explicitly,
\[
 O\!\left(\Lambda^2 L
 [d^3\log^2(2+d)+L(U+1)^2]^2
 (m+1)^4[d+\log(2L)]\right).
\]
The exact ceiling in \eqref{eq:main} includes the harmless lower-order
terms. In particular the degree does not grow with $L$.

\section{Interpretation and remaining questions}
The prior conditions are implementational assumptions. A function may
have both supported and unsupported realizations. Every Boolean function
has some supported realization, but the construction proving this can be
exponentially expensive in input dimension. We have not proved that every
width-$C$, depth-$L$ Boolean circuit has a supported implementation with
energy polynomial in $C,d,L$. Thus this paper supplies a polynomial-depth
\emph{converse}; it does not claim a two-sided polynomial-depth circuit
prior for all small Boolean circuits.

The global gain certificate is sufficient, not necessary for stable
predictions. Saturated tanh circuits can be stable near their actual
states while failing a certificate that allows every diagonal slope in
$[0,1]$. The simpler layerwise nonexpansiveness restriction is stronger
still. Understanding useful weaker restrictions with the same closure
properties is a natural next question.

A uniform output cap is likewise substantive. Without it the correction
amplitude can scale like the square root of current energy, turning the
additive recurrence for its square root into a multiplicative one.
The present proof does not remove that obstruction. Finally, direct sums
need spare neurons. A fixed-width converse needs a way to create that
space without repeatedly paying a constant energy factor, while retaining
the prior's support restriction. That step is not supplied here.

\begin{thebibliography}{9}
\bibitem{quartic} ChatGPT-6 Astra, Claude Fable 5 and Opus 5, Kaarel H\"anni,
and Dmitry Vaintrob. \emph{A quartic-width learning bound for Bayesian
networks stable under internal noise}. Research manuscript, 2026.
\url{https://github.com/kaarelh/AI-works/tree/main/neural-network-generalization/works-in-progress/improved-polynomials/original-support}.
\bibitem{previous} ChatGPT-6 Astra, Kaarel H\"anni, and Dmitry Vaintrob.
\emph{What a few online mistakes imply for a strong-prior neural network}.
Research note, 2026.
\url{https://github.com/kaarelh/AI-works/blob/main/neural-network-generalization/works-in-progress/online-mistakes/nn-converse.md}.
\bibitem{tarres-yao} Pierre Tarr\`es and Yuan Yao.
\emph{Online learning as stochastic approximation of regularization paths:
optimality and almost-sure convergence}.
Appendix A, Corollary A.2, equation (A-1).
\url{https://arxiv.org/abs/1103.5538}.
\end{thebibliography}
\end{document}
